\Article{14}
Denken wir uns das Volumen der beweglichen Nadel in unendlich kleine Elemente zerlegt. Die Koordinaten eines beliebigen Elements seien $x,y,z$, und $e$ sei das darin enthaltene Element des freien Magnetismus. Auf der vertikalen Geraden durch $G$ legen wir den Koordinatenursprung in einen beliebigen Punkt $h$ innerhalb der Nadel. Die Koordinatenachsen $x,y$ sollen horizontal liegen, die erste im magnetischen Meridian nach Norden und die zweite nach Osten gerichtet. Die Koordinate $z$ zählen wir nach oben. Die Wirkung des Erdmagnetismus auf das Element $e$ liefert so zu $\mathrm{d}\Omega$ den Beitrag $Te\mathrm{d}x$.

Ebenso sei das Volumen der zweiten, festen Nadel in unendlich kleine Elemente zerlegt. Einem beliebigen Element entsprechen die Koordinaten $X,Y,Z$ und die Menge an freiem Magnetismus $E$. Schließlich sei $r=\BareRoot{((X-x)^2+(Y-y)^2+(Z-z)^2)}$. Dann liefert die Wirkung des Elements $E$ auf das Element $e$ zur Summe $\mathrm{d}\Omega$ den Beitrag $\frac{eE\mathrm{d}r}{r^n}$, sofern man annimmt, dass diese Wirkung umgekehrt proportional zur Potenz $r^n$ des Abstands $r$ ist.

Bezeichnen wir mit $N$ den Wert von $u$, bei dem der Faden unverdreht ist, so lässt sich das Moment der Torsionskraft des Fadens durch $\theta(N-u)$ ausdrücken. Man kann sich diese Kraft so vorstellen, als wirkte an jedem der beiden Endpunkte des horizontalen Fadendurchmessers am Punkt $G$ eine Tangentialkraft $=\frac{\theta(N-u)}{D}$, wobei $D$ diesen Durchmesser bezeichnet. Daraus ist leicht ersichtlich, dass sich für die Summe $\mathrm{d}\Omega$ der Beitrag $\theta(N-u)\mathrm{d}u$ ergibt.

Die Schwerkraft, die auf die Teilchen der Nadel wirkt, trägt offenbar nichts zur Summe $\mathrm{d}\Omega$ bei, da $u$ die einzige Variable ist. Wir haben daher
\[
\mathrm{d}\Omega=\Sigma Te\mathrm{d}x+\Sigma\frac{eE\mathrm{d}r}{r^n}+\theta(N-u)\mathrm{d}u
\]
wobei im ersten Term über alle Elemente $e$ und im zweiten über alle Paarungen der einzelnen $e$ mit den einzelnen $E$ summiert wird. Die Bedingung für stabiles Gleichgewicht besteht also offenbar darin, dass
\[
\Omega=\Sigma Tex-\Sigma\frac{eE}{(n-1)r^{n-1}}-\tfrac12\theta(N-u)^2
\]
ein Maximum annimmt.

\SourceLeaf{106}
\Article{15}
Für unseren Zweck empfiehlt es sich, die Versuche stets so einzurichten, dass die magnetische Achse jeder der beiden Nadeln horizontal liegt und beide Nadeln sich ungefähr auf derselben Höhe befinden. Auf diese Voraussetzungen beschränken wir daher die weiteren Rechnungen.

Beziehen wir die Koordinaten der Punkte der ersten Nadel auf Achsen, die in dieser Nadel fest sind und sich weiterhin im Punkt $h$ schneiden. Die erste Achse verlaufe in Richtung der magnetischen Achse, die zweite liege horizontal rechts von der ersten und die dritte sei senkrecht nach oben gerichtet. Das Element $e$ habe in Bezug auf diese Achsen die Koordinaten $a,b,c$. Ebenso seien $A,B,C$ die Koordinaten des Elements $E$ in Bezug auf entsprechende Achsen, die in der zweiten Nadel fest sind und sich im Punkt $H$ dieser Nadel schneiden. Wir nehmen an, dass dieser Punkt nahe der Mitte der Nadel und auf derselben Höhe wie der Punkt $h$ gewählt ist.

Ginge es nur um einen Versuch, so ließe sich die Lage des Punkts $H$ am bequemsten durch den Abstand vom Punkt $h$ und die Richtung der Verbindungsgeraden bestimmen. Für unseren Zweck sind jedoch stets mehrere Versuche zu verschiedenen Lagen des Punkts $H$ nötig. Diese Lagen liegen alle auf derselben Geraden, die aber nicht genau durch den Punkt $h$ gehen muss. Daher empfiehlt es sich, die Bezeichnungen von Anfang an so zu wählen, dass eine Reihe solcher Versuche von einer einzigen Variablen abhängt. Wir beziehen daher den Punkt $H$ auf einen beliebigen Punkt $h'$ in derselben horizontalen Ebene nahe bei $h$ mit den Koordinaten $\alpha,\beta,0$ und setzen den Abstand $h'H=R$. Den Winkel zwischen der Geraden $h'H$ und dem magnetischen Meridian setzen wir $=\psi$. Bezeichnen wir nun den Winkel zwischen der magnetischen Achse der zweiten Nadel und dem magnetischen Meridian mit $U$, so erhalten wir
\[
\begin{aligned}
x&=a\cos u-b\sin u\\
y&=a\sin u+b\cos u\\
z&=c\\
X&=\alpha+R\cos\psi+A\cos U-B\sin U\\
Y&=\beta+R\sin\psi+A\sin U+B\cos U\\
Z&=C
\end{aligned}
\]

Damit ist alles zur Entwicklung der Summe $\Omega$ und des Quotienten $\frac{\mathrm{d}\Omega}{\mathrm{d}u}$ vorbereitet, der im Gleichgewicht verschwinden muss.

